\section{A resolution converse for table-dependent bounded errors}
\label{sec:calibration-floor}
The sufficient calibration scale in \eqref{eq:precision} has a converse
in the same bounded-error model. The converse uses physical start
perturbations, not corrupted labels added to an abstract membership
oracle. It concerns short collision commands; arbitrarily long flights
or additional observations of the actual launch position are not included.

\subsection{A common response produced by nearby physical tables}
For a table $\mathcal O_j$ denote its attempted hit bit by $B^j_a(q)$.
An admissible position-error implementation is a measurable map
$q\mapsto F_j(q,a)$ with $|F_j(q,a)-q|\le r$; time and direction
can remain exact. Such maps may depend on the table and the nominal
command. The controller sees only the resulting bit, not the actual
start. This is a submodel of the conditionally bounded, not necessarily
mean-zero couplings in Section~\ref{sec:setting}.

\begin{lemma}[Pointwise common response under two implementations]\label{lem:common-response}
Suppose $\mathcal O_0,\mathcal O_1$ consist of matched strictly convex
components of diameter at most $D_0$, with inradius at least $r_*>0$,
and that matched bodies have Hausdorff distance at most $e>0$.
In each table distinct bodies are separated by at least $d_0$.
Fix $t_{\max}<d_0$. If $e$ is sufficiently small compared with
$r_*$ and $d_0-t_{\max}$, there are measurable position perturbations
of size at most $2e$ such that, for every $|a|\le t_{\max}$ and
with closed-solid and swept-set conventions,
\begin{equation}\label{eq:common-response}
 B^0_a(F_0(q,a))=B^1_a(F_1(q,a))
                    =\min\{B^0_a(q),B^1_a(q)\}.
\end{equation}
The closed convention assigns zero on a solid boundary and one on a
swept boundary outside the solid. In every disagreement case the
forced-zero position is moved strictly into a solid or strictly outside
all sweeps before its response is evaluated. Agreement cases keep the
nominal point and use the same closed convention. The conclusion
applies to locally finite periodic arrays and to both the forward and
translated reverse commands.
\end{lemma}
\begin{proof}
For one fixed $a$ put $S_j(C)=C+[-a,0]$. Matched swept bodies have
Hausdorff distance at most $e$, because their support functions differ
by the same amount as those of the matched bodies. In either table,
\begin{equation}\label{eq:swept-separation}
 \dist(S_j(C),S_j(D))\ge d_0-|a|\quad(C\ne D).
\end{equation}
Indeed two points in the sweeps are $c-ua$ and $d-va$, with
$u,v\in[0,1]$. Their distance is at least $|c-d|-|u-v||a|$.
Thus a start belonging to a swept body determines that body uniquely.
Moreover no solid other than that body meets this sweep. On it the
hit bit is the indicator of the sweep minus the indicator of its body.

Leave the start unchanged when the two bits agree, and also on the
side whose bit is already zero. Consider $B^0_a(q)=1$, $B^1_a(q)=0$,
and let $C_0$ be the unique body whose sweep contains $q$. Write $C_1$
for its matched body. The separation \eqref{eq:swept-separation} and
Hausdorff matching imply the following exhaustive alternatives:
\[
                   q\in C_1,\qquad\hbox{or}\qquad q\notin\operatorname{int}S_1(C_1).
\]
To check this, every other sweep in table 1 is at distance at least
$d_0-|a|-e$ from $S_0(C_0)$, so cannot supply either a competing hit
or a solid start there. If $q\in S_1(C_1)\setminus C_1$, its bit
would be one.

In the first case $\dist(q,C_0)\le e$. Project $q$ onto $C_0$, and
move the projection a further distance less than $e/2$ toward a fixed
interior point of $C_0$. This produces an interior solid start within
$3e/2$ of $q$, whose bit is zero. Uniform inballs permit a fixed small
convex-combination factor. The construction is measurable.

In the second case, when $q$ is outside the closed sweep, let $q_1$
be its metric projection onto $S_1(C_1)$ and set
$v=(q-q_1)/|q-q_1|$. At a boundary point choose an outward unit normal
from its compact normal cone, measurably, for example by the least
angle in a fixed angular convention. In both cases
$q\cdot v\ge p_{S_1(C_1)}(v)$, and support matching gives
$p_{S_0(C_0)}(v)\le p_{S_1(C_1)}(v)+e$. Hence
\[
                       q'=q+\tfrac32 e v
\]
lies strictly outside $S_0(C_0)$. It cannot enter another swept body
of table 0, by \eqref{eq:swept-separation} and smallness of $e$.
It is therefore a free miss, with bit zero. The displacement is $3e/2$.
Projection onto a closed convex set is continuous, so this map is
measurable on its case domain. Reverse the indices when the differing
bit is on side 1. Local finiteness and the separation give measurable
component selection. For $a=0$ every bit is zero and no change is needed.
Boundary cases can be assigned by the same inward or outward moves;
in particular the final forced-zero positions are not on boundaries.

The proof is pointwise in the nominal start and displacement. Apply it
to $(q,a)$ for a forward command and to $(q+a,-a)$ for its reverse.
Fresh nominal draws and deterministic perturbation maps preserve
conditional independence. They do not supply another observable.
\end{proof}

\begin{theorem}[Minimax resolution under table-dependent bounded errors]
\label{thm:calibration-floor}
For each fixed $s=6+\beta$ there is a fixed physical subclass of
$\mathcal B_\eta$, with known primitive lattice, two known species
and the fixed laboratory frame, and constants $c,r_0>0$ such that
for every $0<r<r_0$ two tables in this class have centered-support
$C^2$ distance at least
\begin{equation}\label{eq:calibration-separation}
                            c r^{(s-2)/s}
\end{equation}
yet admit position-error implementations bounded by $r$ with identical
laws for the entire adaptive bit experiment. This holds for every
nominal command design with $|a|\le t_{\max}<d_0$, arbitrary input
precision and arbitrary numbers of attempts. A sequential stopping
rule based on the bits cannot distinguish them either.

To specify the uncertainty quantifiers, let $\mathfrak I_r(\mathcal O)$
be the measurable implementation maps permitted above, and let
$\mathcal A_{t_{\max}}$ be the admissible bit-based controllers. The
construction has the stronger, controller-uniform form
\begin{gather}\label{eq:calibration-quantifiers}
 \forall\,0<r<r_0\quad
 \exists\,(\mathcal O_0,F_0),(\mathcal O_1,F_1),\quad
             F_j\in\mathfrak I_r(\mathcal O_j),\notag\\
 \forall\,\mathcal A\in\mathcal A_{t_{\max}}:\qquad
 \mathsf{Law}_{\mathcal O_0,F_0}^{\mathcal A}(Z)
       =\mathsf{Law}_{\mathcal O_1,F_1}^{\mathcal A}(Z).
\end{gather}
Here $Z$ includes the returned bits and the controller's recorded
commands and stopping decisions, not the actual launch positions.
The maps $F_0,F_1$ need not agree and may use the nominal draw and
command. They are fixed measurable maps for the two worlds, not an
extra observation or a common known jitter distribution.

Consequently, uniform success probability greater than $1/2$ at
centered $C^2$ accuracy $\nu$ under all such position errors requires
\begin{equation}\label{eq:calibration-necessary}
                            r\le C\nu^{s/(s-2)}.
\end{equation}
Constants, not an equality threshold, are asserted. The result is a
worst-case bounded-calibration statement; it does not assert this
obstruction for one fixed common offset or for a known mean-zero law.
\end{theorem}
\begin{proof}
Use the two-species periodic construction in Lemma~\ref{lem:packing},
but compare the unit-disk support $p_0=1$ with a single nonzero bump
\[
 p_1(\theta)=1+a h^s\psi((\theta-\theta_0)/h).
\]
The amplitude $a$ is fixed small enough that all curvature, separation
and patch margins remain uniform. Choose
$h=(r/(8a\|\psi\|_\infty))^{1/s}$. Then the Hausdorff distance of
the two variable bodies is at most $e=r/8$. The fixed disk and all
periodic copies are matched without alteration. At the bump center
the second derivatives differ by $a|\psi''(0)|h^{s-2}$; the
first-harmonic estimate in Lemma~\ref{lem:centered-packing} shows that
centering removes only $O(h^s)$. This gives
\eqref{eq:calibration-separation}. The lattice and disclosed data are
identical for both tables; the bump does not change primitiveness.

Lemma~\ref{lem:common-response} produces two physical implementations
with position errors at most $2e<r$. Couple the parameter-independent
controller randomness and every nominal draw across the two worlds.
Their first bits are equal by \eqref{eq:common-response}. Inductively,
identical histories generate identical next commands and identical
bits. Thus the complete bit streams, nominal controller records and
all bit-measurable stopping decisions have identical laws. No finite
or infinite sequence of those observations distinguishes the worlds.
An output law common to two disjoint accuracy sets cannot place
probability greater than $1/2$ in each. If $2\nu$ is smaller than
\eqref{eq:calibration-separation}, uniform success greater than $1/2$
is impossible. Rearranging proves \eqref{eq:calibration-necessary}.
\end{proof}

\subsection{The two resource requirements}
Together with Theorems~\ref{thm:adaptive} and \ref{thm:expected}, this
gives two necessary resources on one fixed nondegenerate subclass:
\[
 \sup_{\mathcal O}\E_{\mathcal O}\mathsf T
       \ge c\nu^{-1/(s-2)},\qquad
 r\le C\nu^{s/(s-2)}.
\]
Conversely the reciprocal compass construction attains the geometric
accuracy with the stated logarithmic overhead in attempts when
$\ell+\tau+t\alpha\le c\nu^{s/(s-2)}$. The converse calibration
necessity already holds with duration and angle exact, so it concerns
position error alone, not separate optimal exponents for all three
actuators. The lower bound allows any directions of length at most
$t_{\max}$, a larger design class than the four-atom upper construction.
The finite digital realization in Section~\ref{sec:digital} does not
alter either bound.

Thus the shrinking physical tolerance is not merely an omitted term in
a sample count: under the stipulated worst-case coupling it cannot be
replaced by additional samples. This is a statement within the active
collision model. Periodicity, the bounded presentation and the known
patch margin are still priors for the uniform global conclusion.
Neither crystallinity, apparatus calibration nor passive spectral
rigidity is inferred from the returned bits.
